\begin{enumerate}[label=\textbf{\ansnum.\arabic*}]
  \item \verb!Option<T>! says ``there may be no value'' without a reason: the ball is not visible, a list is empty. \verb!Result<T, E>! says ``this operation can fail, and here is why'': a packet that does not parse (the error says which field), a configuration file that cannot be opened (the error says that the file does not exist or may not be read). Rule: if the caller would want to know the reason, use \verb!Result!.
  \item
%rs:a2
\begin{verbatim}
use std::num::ParseFloatError;

fn parse_pair(s: &str) -> Result<(f64, f64), ParseFloatError> {
    let mut it = s.split(',');
    let a = it.next().unwrap_or("").trim().parse()?;
    let b = it.next().unwrap_or("").trim().parse()?;
    Ok((a, b))
}

fn main() {
    println!("{:?}", parse_pair("1.5,-2"));
    println!("{:?}", parse_pair("1.5"));
    println!("{:?}", parse_pair("x,1"));
}
\end{verbatim}
prints \verb!Ok((1.5, -2.0))! and then two times \verb!Err(ParseFloatError { kind: ... })! (\verb!Empty! for the missing part of \verb!"1.5"!, \verb!Invalid! for \verb!"x"!). A missing part is treated as the empty text, whose parsing fails. (A better error type would say ``missing second value''; see the enum of the chapter.)
  \item
%rs:a3
\begin{verbatim}
fn main() {
    let d = vec![0.3, 0.8, 1.2, 2.5, 1.9];
    let n = d.iter().filter(|&&x| (0.5..=2.0).contains(&x)).count();
    let above: Vec<f64> = d.iter().copied().filter(|&x| x > 1.0).collect();
    let mean = above.iter().sum::<f64>() / above.len() as f64;
    let squares: Vec<f64> = d.iter().map(|x| x * x).collect();
    println!("{n} {mean:.3} {squares:?}");
}
\end{verbatim}
prints $3$, $1.867$ and the squares $[0.09, 0.64, 1.44, 6.25, 3.61]$ (up to floating-point rounding in the last digits). The values in $[0.5,2.0]$ are $0.8$, $1.2$ and $1.9$.
  \item
%rs:a4
\begin{verbatim}
fn dot(x: &[f64], y: &[f64]) -> f64 {
    x.iter().zip(y).map(|(a, b)| a * b).sum()
}

fn main() {
    println!("{}", dot(&[1.0, 2.0, 3.0], &[4.0, 5.0, 6.0]));
}
\end{verbatim}
The result is $1\cdot4+2\cdot5+3\cdot6=32$. If the vectors have different lengths, \verb!zip! silently stops at the shorter one: add an \verb|assert_eq!(x.len(), y.len())| if that is not what you want.
  \item
%rs:a5
\begin{verbatim}
fn moving_average(xs: &[f64], n: usize) -> Vec<f64> {
    xs.windows(n).map(|w| w.iter().sum::<f64>() / n as f64).collect()
}

fn main() {
    println!("{:?}", moving_average(&[1.0, 2.0, 3.0, 4.0, 5.0], 3));
}
\end{verbatim}
The result is $[2,3,4]$ (three values: $5-3+1$). \verb!windows(0)! panics, so a robust version checks \verb!n > 0!.
  \item
%rs:a6
\begin{verbatim}
use std::collections::HashMap;

fn main() {
    let text = "the ball the goal the robot a ball";
    let mut count: HashMap<&str, usize> = HashMap::new();
    for w in text.split_whitespace() {
        *count.entry(w).or_insert(0) += 1;
    }
    let mut v: Vec<(&str, usize)> = count.into_iter().collect();
    v.sort_by(|a, b| b.1.cmp(&a.1).then(a.0.cmp(b.0)));
    for (w, n) in v { println!("{w} {n}"); }
}
\end{verbatim}
prints \verb!the 3!, \verb!ball 2!, \verb!a 1!, \verb!goal 1!, \verb!robot 1!. The comparison first looks at the counts in reverse order, and \verb!then! (only for equal counts) at the words.
  \item (a) Result (bad data from outside is normal). (b) Panic: if the index can be $\ge22$ it is a bug in the program, and a stop with a clear message is better than wrong data; an \verb!assert! or the normal index check does that. (c) Result (the user may have misspelled the name). (d) Panic: a violated, documented precondition is a bug in the caller (one may also return an \verb!Option!, if the caller cannot know). (e) Result or \verb!Option! (cameras fail; the robot must go on, e.g.\ with the last image or a fallback).
  \item
%rslib:a7
\begin{verbatim}
/// Limits `x` to `lo..=hi`. Panics if `lo > hi`.
pub fn clamp(x: f64, lo: f64, hi: f64) -> f64 {
    assert!(lo <= hi, "lo must not be larger than hi");
    x.max(lo).min(hi)
}

#[cfg(test)]
mod tests {
    use super::*;
    #[test]
    fn inside() { assert_eq!(clamp(0.3, 0.0, 1.0), 0.3); }
    #[test]
    fn below() { assert_eq!(clamp(-2.0, 0.0, 1.0), 0.0); }
    #[test]
    fn above() { assert_eq!(clamp(7.0, 0.0, 1.0), 1.0); }
    #[test]
    #[should_panic(expected = "lo must not")]
    fn reversed_limits() { clamp(0.5, 1.0, 0.0); }
}
\end{verbatim}
All four tests pass. The decision (panic on reversed limits) is a programming error of the caller and is documented in the \verb!///! comment, which \verb!cargo doc! shows.
  \item \verb!Cargo.toml! of the workspace: \verb![workspace]!, \verb!members = ["geometry", "app"]!, \verb!resolver = "3"!. In \verb!geometry/Cargo.toml! only the \verb![package]! section (name, version, edition) and the code in \verb!geometry/src/lib.rs! (\verb!pub fn distance(...)!). In \verb!app/Cargo.toml!, below \verb![package]!: \verb![dependencies]! and \verb!geometry = { path = "../geometry" }!. In \verb!app/src/main.rs!: \verb!geometry::distance(...)!. This was built and run for the book; \verb!cargo run -p app! printed \verb!5! for the points $(0,0)$ and $(3,4)$, and \verb!cargo build --release! builds both crates.
  \item[\textbf{\ansnum.10}$\star$]
%rslib:a10
\begin{verbatim}
use std::fmt;

#[derive(Debug, PartialEq)]
pub enum CsvError { BadLine(usize), NoHeader }

impl fmt::Display for CsvError {
    fn fmt(&self, f: &mut fmt::Formatter) -> fmt::Result {
        match self {
            CsvError::BadLine(n) => write!(f, "line {n}: expected 'time,x'"),
            CsvError::NoHeader => write!(f, "file has no header line"),
        }
    }
}

pub fn read_csv(text: &str) -> Result<Vec<(f64, f64)>, CsvError> {
    let mut lines = text.lines().enumerate();
    lines.next().ok_or(CsvError::NoHeader)?;          // skip the header
    let mut out = Vec::new();
    for (i, line) in lines {
        if line.trim().is_empty() { continue; }
        let mut p = line.split(',');
        let (a, b) = (p.next(), p.next());
        match (a.and_then(|s| s.trim().parse().ok()),
               b.and_then(|s| s.trim().parse().ok()), p.next()) {
            (Some(t), Some(x), None) => out.push((t, x)),
            _ => return Err(CsvError::BadLine(i + 1)),
        }
    }
    Ok(out)
}

#[cfg(test)]
mod tests {
    use super::*;
    #[test]
    fn good_file() {
        let r = read_csv("time,x\n0,1.5\n\n0.1,2\n").unwrap();
        assert_eq!(r, vec![(0.0, 1.5), (0.1, 2.0)]);
    }
    #[test]
    fn text_instead_of_number() {
        let r = read_csv("time,x\n0,1\n1,abc\n");
        assert_eq!(r, Err(CsvError::BadLine(3)));
    }
    #[test]
    fn empty_file() {
        assert_eq!(read_csv(""), Err(CsvError::NoHeader));
    }
}
\end{verbatim}
All three tests pass. The line number counts from $1$ including the header, so the error message matches what the user sees in an editor. The three-way \verb!match! on a tuple checks that there are exactly two parseable numbers and nothing more.
\end{enumerate}
